\documentclass[a4paper, 12pt]{article}

\usepackage[utf8]{inputenc}
\usepackage[T1]{fontenc}
\usepackage{textcomp}
\usepackage{amssymb}
\usepackage{newtxtext} \usepackage{newtxmath}
\usepackage{amsmath, amssymb}
\newtheorem{problem}{Problem}
\newtheorem{example}{Example}
\newtheorem{lemma}{Lemma}
\newtheorem{theorem}{Theorem}
\newtheorem{problem}{Problem}
\newtheorem{example}{Example} \newtheorem{definition}{Definition}
\newtheorem{lemma}{Lemma}
\newtheorem{theorem}{Theorem}
\DeclareMathAlphabet{\mathcal}{OMS}{cmsy}{m}{n}

\begin{document}

\textit{La no-resistencia al mal.} Refiriéndose a los practicantes de la
no-resistencia al mal, Tolstoi nos dice que incluso si fueran una diminuta
minoría, sujeta a nada más que el ridículo y el desprecio del mundo, el mundo,
«sin notarlo y sin sentirse obligado a ello, se volvería más sabio y mejor por
obra de esta influencia secreta». Jesús, que como mínimo era un hombre sabio,
enseñó que «somos la luz del mundo»: basta con practicar una ética compasiva
para dar luz a las tinieblas. El sabio checo Petr Chelčický, una figura curiosa
y largamente ignorada, enseñó que todo practicante de una ética cristiana
debería negarse a la violencia y, en particular, al servicio militar, notando
que si los pobres se negaran los ricos no tendrían a quién mandar a luchar por
ellos. E incluso cierto guerrero, sin duda ajeno a la doctrina de la
no-resistencia al mal, enseñó que la ternura es el bien que nunca debe perderse. 

~

La no-resistencia al mal es una doctrina poco difundida, todavía menos
comprendida, y frecuentemente ridiculizada. Es, sin embargo, la única ética
superadora de la violencia, el ejercicio del poder, y la dominación de fuertes
sobre débiles. Y cuantos más las practiquemos, más sentirá el mundo la
«influencia secreta» de la compasión y el amor.








\end{document}



